\documentclass[12pt]{article}
\usepackage[utf8]{inputenc}
\usepackage[T1]{fontenc}
\usepackage{hyperref}
\title{Betta Kurumba Language Resources for K-12 Curriculum Development}
\author{Compiled by Claude Sonnet 5.5 (claude-sonnet-5-5) for Andr\'as Kornai}
\date{}
\begin{document}
\maketitle

\section{Introduction}
Betta Kurumba (ISO 639-3 xub, Glottocode bett1235) is a Dravidian language of the Nilgiri region of India, with about 32,000 speakers. Some sources say the community may speak a dialect of Kannada, so its classification is not settled. This survey is based on one web search, and I did not open the linked PDFs or archives.

\section{Primary Educational Resources}
No primers, textbooks or school materials were found.

\section{Linguistic and Reference Materials}
\begin{itemize}
\item Wikipedia, Betta Kurumba language, which points to Gail Coelho's grammar (2003, archived as a PDF at the University of Texas library) and her 2018 book of annotated texts. \url{https://en.wikipedia.org/wiki/Betta_Kurumba_language}
\item Glottolog, Betta Kurumba (genealogy, identifiers, endangerment status). \url{https://glottolog.org/resource/languoid/iso/xub}
\item Endangered Languages Project, Betta Kurumba, which lists a 2012 survey by Gnanasundaram, Perialwar and Rangan in the International Journal of Dravidian Linguistics. \url{https://www.endangeredlanguages.com/lang/1453?hl=fr}
\end{itemize}

\section{Religious and Community Resources}
\begin{itemize}
\item Global Recordings Network, Kurumba, Vatta (Bible-based audio and slideshow materials). \url{https://globalrecordings.net/en/language/xub}
\end{itemize}

\section{Digital and Audio Resources}
\begin{itemize}
\item CIIL Bhasha Sanchika, Betta Kuruba community page. \url{https://sanchika.ciil.org/communities/f471779a-3400-4462-981d-c95b041550b4}
\end{itemize}

\section{Curriculum Development Context}
The language is not known to be taught in schools. The 2012 survey reports that younger speakers could not produce some basic vocabulary, which suggests that language transmission is weakening. Coelho's grammar and texts and the CIIL collection are the most useful starting points. Any curriculum would need to be developed with the community.
\end{document}